% sec7-errata.tex -- Section 7 of the C3 paper (draft 2026-11-08): result G.
% Sources: codex-20261104-runB-adjudication/03-C1-INTERFACES-AND-ERRATA.txt \S5 (wording limits);
% notes-20261104/referee-runB-closure-periodicity.md; claude-20261104-c3-closure/01 \S4 R1, \S6;
% codex-20261104-c3-monoidal/03-BOUNDARY-AND-STRICT-REPAIR.md.

\section{Errata in two sources}\label{sec:errata}

Writing the proofs of Section~\ref{sec:periodicity} required checking two recent sources
closely.  We record five points.  For each we say whether it is a counterexample or an
omitted justification, and we do not draw conclusions beyond the stated point; in
particular none of them suggests that the other results of either paper fail.  The
first four points were sent on 8 November 2026 to Harman and to Yacobi, one of the authors
of \cite{hone2025}.  The fifth was found afterwards, in the reconstruction of
\cite[Thm.~8.4.5]{hone2025} described in Section~\ref{sec:periodicity}.

\paragraph{Harman, Corollary~1.8(4), at $p=2$ (counterexample).}  The corollary identifies
the value of a right adjoint on a Young module with a tilting module,
$\tilde G(Y(\lambda))=T(\lambda')$~\cite[Thm.~1.7(4), Cor.~1.8(4)]{harman2015}.  At $p=2$ in
polynomial degree $2$, $Y(2)$ is the trivial module $k$; the right-adjoint value has
dimension $N(N+1)/2$, while $T(1,1)=\Lambda^2V$ has dimension $N(N-1)/2$ (at $N=2$: $3$ against
$1$).  Harman's proof of his Theorem~2.8 uses this identification to conclude that the
twisted Schur functor is fully faithful on tilting modules.  That conclusion is true in every
characteristic: the tilting modules in question are the summands of the $\Lambda^\lambda V$;
the ordinary Schur functor sends $\Lambda^\lambda V$ to $M(\lambda)\otimes\mathrm{sgn}$ and
the twisted one to $M(\lambda)$; the
Hom spaces on both sides have the dimension of the number of $(\lambda,\mu)$ double cosets;
and the functor is faithful on them because $\Lambda^\mu V\subseteq V^{\otimes d}$ and
$V^{\otimes d}$ is projective and injective with socle detected by the Schur functor.  The
second proof of Section~\ref{sec:periodicity-tilting} uses this argument instead of
Corollary~1.8(4).

\paragraph{Harman, Theorem~2.5, the modulus (counterexample at the printed period).}  The
theorem's printed modulus allows, at $r=p=3$, the period $3$.  But the stable incidence
composite through $3$-subsets has scalar $\binom73\equiv2$ at rank $7$ and
$\binom{10}3\equiv0$ at rank $10$ in $\mathbb F_3$; both ranks exceed $2r$, so no functor
identifying these stable maps can respect this composite at period $3$.  A period that is a
power of $p$ strictly greater than $r$ avoids the problem, and that is what we use.  We do not know whether this has been noted before; we record it because our statement
depends on the correct modulus.

\paragraph{Harman, Theorem~2.9 (omitted justification).}  The proof applies Lemma~2.10 to a
restricted right adjoint without showing that it lands in the Serre subcategory where the
lemma applies.  The conclusion needed, that the Serre quotient by the kernel of the twisted
Schur functor is its full image with ambient kernels and subobjects, follows from the
pullback-and-graph argument of Section~\ref{sec:periodicity-oti} (compare
\cite[Lemma~8.5.1]{hone2025}).  This is not a counterexample to Theorem~2.9.

\paragraph{Hone et al., Theorem~8.5.3, at the boundary of Definition~6.1.1
(counterexample).}  Take $p=q=2$ and $d=4$; the label $(3,1)$ is $2$-stable, since
$\lambda_1=3\ge(4+2)/2$, and it truncates to the $2$-singular label $(1,1)$.  Let $W$ be the
sum-zero subspace of $\mathbb F_2^4$ (a submodule of the permutation module of $S_4$) and $C$
the line of the all-ones vector.  In $W/C$ with basis $a=e_1+e_2$, $b=e_1+e_3$, the
transposition $(34)$ acts by $a\mapsto a$, $b\mapsto b+a$.  So $W/C$ is free over
$\mathbb F_2[\langle(34)\rangle]$ and is killed by semisimplification, although it is a
non-zero object of the stable image category: it is a subquotient of $E_1(4)$, which belongs to
that category, and the category is closed under ambient subquotients.  The equivalence asserted in
\cite[Thm.~8.5.3]{hone2025} is therefore not faithful at that boundary.  The example does not
refute Theorem~8.4.5 of the same paper: the descended functor can be exact while killing this
object.  This is why the proof of Section~\ref{sec:periodicity-oti} detects the constant
through a map out of the simple unit and never uses faithfulness.  With strict inequality in
Definition~6.1.1(2b) the boundary case does not arise.

\paragraph{Hone et al., Lemma~8.2.1 (counterexample), and the proof of Theorem~8.4.5.}
Lemma~8.2.1 of \cite{hone2025} states: for a highest-weight category $(\mathcal C,\Lambda)$, a
maximal $\lambda$, and an additive functor $\varphi$ to an abelian category that is exact on
the Serre subcategory of $\Lambda\setminus\{\lambda\}$ and whose comparison maps
$\varphi(\Delta_\lambda)\to\tilde\varphi(\Delta_\lambda)$ and
$\tilde\varphi(\nabla_\lambda)\to\varphi(\nabla_\lambda)$ are isomorphisms ($\tilde\varphi$
the tilting derived functor), the functors $L\varphi$, $\tilde\varphi$, $R\varphi$ agree and
$\tilde\varphi$ is a direct summand of $\varphi$.  This is false.  Let $\mathcal C$ be the
representations of the Kronecker quiver (vertices $1,2$, two arrows $1\to2$) over any field,
ordered $1<2$.  Then $\Delta_1=\nabla_1=S_1$, $\Delta_2=S_2$ is projective, $\nabla_2=I_2$ is
the injective of dimension vector $(2,1)$, and the tilting objects are the sums of $S_1$ and
$I_2$.  Let $\psi(X)$ be the sum of the images of all maps from $S_1\oplus S_2\oplus I_2$ to
$X$.  It is additive, exact on the Serre subcategory generated by $S_1$, and the identity on
$S_1$, $S_2$, $I_2$ and the maps between them, so both hypotheses hold and $\tilde\psi$ is the
identity.  For the projective cover $P_1$ of $S_1$, of dimension vector $(1,2)$,
$\mathrm{Hom}(S_1,P_1)=\mathrm{Hom}(I_2,P_1)=0$ and $\mathrm{Hom}(S_2,P_1)=k^2$, so
$\psi(P_1)=S_2^2$, of which $P_1=\tilde\psi(P_1)$ is not a direct summand; and
$L\psi(S_1)=0\ne S_1$.  The error in the proof is the claim that $P_\bullet\to
i_*i^*P_\bullet$ is a homotopy equivalence: a short exact sequence of complexes with
contractible kernel gives a quasi-isomorphism, which an additive functor need not preserve.
Lemma~8.4.4 uses Lemma~8.2.1 to identify the OTI functor with its tilting derived functor on
all stable objects, and Theorem~8.4.5 rests on Lemma~8.4.4.  So the published proof of
Theorem~8.4.5 is incomplete.  This is not a counterexample to Theorem~8.4.5, which may hold
for the OTI functor by an argument specific to it.  The point was found by Astra in a model
reconstruction of the proof and checked by hand by a second model.

Further minor points found in the same reading (a misprint in \cite[\S7.3]{hone2025},
where $M^{\lambda^t}\otimes\mathrm{sgn}$ should read $M^\lambda\otimes\mathrm{sgn}$; the
justification of \cite[Prop.~2.7]{harman2015} at $\ell=1$) and, in the same reconstruction, two points
in \cite[\S8.3, \S7.3]{hone2025} (the genericity condition before Definition~8.3.1 is not
sufficient for Lemma~8.3.2 in characteristic $2$, while coefficients linearly independent
over $\mathbb F_p$, which the paper gives as an example, suffice; and in the proof of
Lemma~7.3.4 the dominance inequality should read $\ge$, with target degree $M-q$) do not
affect anything here and were not included in those notes.
